\chapter{Connectivity and Graph Features}
\label{chap:connectivity-graph-features}

\section{Overview of Network Construction}
\label{sec:functional-network-construction-overview}

After preprocessing and temporal segmentation, each seizure-related EEG window
was represented as a three-dimensional array containing \num{90}
non-overlapping epochs, \num{22} bipolar EEG channels, and \num{512} samples per
epoch. These segmented EEG windows formed the input to the functional
connectivity pipeline. The purpose of this stage was to transform each
time-windowed EEG segment into a set of functional brain networks that could be
summarized using graph-theoretical measures.

In this thesis, a functional brain network is defined as a weighted graph
constructed from pairwise statistical relationships between EEG channels. Each
of the \num{22} bipolar channels corresponds to one graph node. The edges
between nodes represent functional connectivity values estimated between pairs
of channels. Therefore, for each seizure window, frequency band, and
connectivity method, the pipeline produced one square connectivity matrix of
size \num{22} $\times$ \num{22}. Each entry of this matrix represents the
estimated functional coupling between two EEG channels.

Connectivity was computed separately for each temporal window, frequency band,
and connectivity method. The four temporal windows were Baseline, T0, T1, and
T2, as defined in \cref{sec:temporal-window-definition}. The six analyzed
frequency bands were delta, theta, alpha1, alpha2, beta1, and beta2, as defined
in \cref{sec:frequency-band-decomposition}. For each saved EEG segment, three
connectivity methods were applied: magnitude-squared coherence, weighted
phase-lag index, and amplitude envelope correlation. Consequently, each saved
window produced 18 functional connectivity matrices.

The complete network construction process can therefore be summarized as
follows. First, each segmented EEG window was decomposed according to the
predefined frequency bands. Second, functional connectivity was estimated
between all pairs of channels using each connectivity method. Third, the
resulting connectivity matrices were validated to ensure that they had the
expected \num{22} $\times$ \num{22} shape, contained only finite values, and were
symmetric. Finally, each matrix was interpreted as a weighted undirected network
and passed to the graph feature extraction stage.

This organization allowed the analysis to examine whether network properties
changed across time windows before seizure onset, whether these changes were
frequency-specific, and whether they depended on the connectivity method used
to define the network. It also preserved the link between each graph, the
corresponding subject/case identifier, seizure, temporal window, frequency
band, and connectivity method, which was essential for the later statistical
analysis.

\section{Functional Connectivity Estimation}
\label{sec:functional-connectivity-estimation}

Functional connectivity was estimated between all pairs of EEG channels for
each saved temporal window and frequency band. Three complementary
connectivity measures were used: magnitude-squared coherence, weighted
phase-lag index (wPLI), and amplitude envelope correlation (AEC). These
methods were selected because they capture different aspects of statistical
dependence between EEG signals: coherence reflects frequency-specific linear
coupling, wPLI emphasizes consistent non-zero phase-lag synchronization, and
AEC captures co-fluctuations in band-limited amplitude envelopes
\cite{Friston1994,Vinck2011,Bruns2000}.

The use of connectivity-derived networks is consistent with previous epilepsy
studies that model EEG or intracranial EEG activity as functional graphs, where
nodes represent recording sites and edges represent pairwise functional
coupling \cite{Vecchio2016,Lanzone2022}.

For each segmented EEG window, connectivity was computed separately for the
six predefined frequency bands described in
\cref{sec:frequency-band-decomposition}. The output of each method was a
symmetric \num{22} $\times$ \num{22} matrix, where rows and columns corresponded
to the fixed bipolar EEG channels. The diagonal entries were set to one before
saving, since each channel is perfectly connected with itself by definition;
self-connections were later removed before graph-theoretical analysis.

\subsection{Coherence}
\label{subsec:coherence}

Magnitude-squared coherence was used to estimate frequency-specific linear
coupling between pairs of EEG channels \cite{Pfurtscheller1999,Vecchio2016}. For two signals $x_i(t)$ and $x_j(t)$,
coherence at frequency $f$ can be expressed as

\[
    C_{ij}(f) =
    \frac{|S_{ij}(f)|^2}{S_{ii}(f)S_{jj}(f)},
\]

where $S_{ij}(f)$ is the cross-spectral density between channels $i$ and $j$,
and $S_{ii}(f)$ and $S_{jj}(f)$ are the corresponding auto-spectral densities.
The value of coherence lies between zero and one, with higher values indicating
stronger linear coupling at the given frequency.

In the implemented pipeline, cross-spectral and auto-spectral densities were
estimated using Welch-based spectral estimation. For each \SI{2}{\second}
epoch, the spectral estimates were computed using a segment length of
\num{256} samples and an overlap of \num{128} samples. Coherence values were
then averaged across all epochs and across all frequency bins belonging to the
selected frequency band. This produced one coherence value for each channel
pair, frequency band, and temporal window.

The resulting coherence matrices were clipped to the interval $[0,1]$ and
validated before being saved. Because coherence is a symmetric pairwise measure,
the resulting matrices were symmetric by construction and represented weighted
undirected functional networks.

\subsection{Weighted Phase-Lag Index}
\label{subsec:wpli}

The weighted phase-lag index was used to estimate phase-based synchronization
while reducing the influence of zero-lag coupling \cite{Vinck2011,Lanzone2022}. This is important in scalp
EEG, where instantaneous coupling can arise from volume conduction, common
reference effects, or shared sources. Unlike coherence, which can be influenced
by both zero-lag and non-zero-lag relationships, wPLI emphasizes consistent
phase differences with a non-zero imaginary component.

Before computing wPLI, the EEG epochs were band-pass filtered within the
frequency band of interest using a fourth-order Butterworth filter. The analytic
signal was then obtained using the Hilbert transform. For a pair of analytic
signals, the cross-spectrum-like quantity was computed and its imaginary part
was used to estimate wPLI:

\[
    \mathrm{wPLI}_{ij}
    =
    \frac{
        \left|
        \mathbb{E}\left[\operatorname{Im}(Z_{ij})\right]
        \right|
    }{
        \mathbb{E}\left[
        \left|\operatorname{Im}(Z_{ij})\right|
        \right]
    },
\]

where $Z_{ij}$ denotes the product of the analytic signal from channel $i$ and
the complex conjugate of the analytic signal from channel $j$, and
$\operatorname{Im}(\cdot)$ denotes the imaginary part. The expectation was
computed across epochs and time samples.

The resulting wPLI values range from zero to one. Values close to zero indicate
weak or inconsistent phase-lagged coupling, whereas values closer to one
indicate more consistent non-zero phase-lag synchronization between two
channels. As with coherence, the output was a symmetric \num{22} $\times$
\num{22} connectivity matrix for each temporal window and frequency band.

\subsection{Amplitude Envelope Correlation}
\label{subsec:aec}

Amplitude envelope correlation was used to estimate coupling between
band-limited amplitude fluctuations. While coherence and wPLI focus on spectral
or phase-based relationships, AEC captures whether the amplitude envelopes of
two EEG channels fluctuate together over time within a given frequency band \cite{Bruns2000,Hipp2012}.

For each frequency band, the EEG epochs were first band-pass filtered using the
same fourth-order Butterworth filtering procedure used for wPLI. The analytic
signal was then computed with the Hilbert transform, and the amplitude envelope
of each channel was obtained by taking the absolute value of the analytic
signal. For channel $i$, the amplitude envelope can be written as

\[
    A_i(t) = |z_i(t)|,
\]

where $z_i(t)$ is the analytic signal of the band-limited EEG signal.

For each epoch, Pearson correlation coefficients were computed between the
amplitude envelopes of all channel pairs. The final AEC matrix for a given
window and frequency band was obtained by averaging these epoch-level
correlation matrices. Unlike coherence and wPLI, AEC values may be positive or
negative, because they are based on correlation coefficients. Therefore, the
values were clipped to the interval $[-1,1]$ before validation and saving.

AEC provides a complementary view of functional connectivity because it is
sensitive to co-fluctuations in signal amplitude rather than phase consistency
alone. However, because no orthogonalization or source reconstruction was
applied, AEC may still be influenced by zero-lag correlations and shared
sensor-level activity. This limitation is considered when interpreting
connectivity differences across temporal windows.

\section{Adjacency Matrix Construction}
\label{sec:adjacency-matrix-construction}

After functional connectivity estimation, each temporal window, frequency band,
and connectivity method was represented by a square adjacency matrix. These
matrices provided the bridge between EEG signal analysis and graph-theoretical
feature extraction. In this representation, EEG channels were treated as graph
nodes and pairwise connectivity estimates were treated as weighted edges.

For each saved EEG segment, the pipeline produced one adjacency matrix for each
combination of frequency band and connectivity method. Since the analysis used
\num{6} frequency bands and \num{3} connectivity methods, each temporal window
produced \num{18} adjacency matrices. Each matrix had dimensions
\num{22} $\times$ \num{22}, corresponding to the fixed \num{22}-channel bipolar
montage described in \cref{subsec:channel-selection-montage-standardization}.

\subsection{Weighted Matrices}
\label{subsec:weighted-matrices}

The connectivity matrices were treated as weighted adjacency matrices. For a
given seizure, temporal window, frequency band, and connectivity method, the
matrix can be written as

\[
    \mathbf{A} =
    \begin{bmatrix}
    A_{11} & A_{12} & \cdots & A_{1N} \\
    A_{21} & A_{22} & \cdots & A_{2N} \\
    \vdots & \vdots & \ddots & \vdots \\
    A_{N1} & A_{N2} & \cdots & A_{NN}
    \end{bmatrix},
\]

where $N = 22$ is the number of EEG channels. Each off-diagonal element
$A_{ij}$ represents the estimated functional connectivity between channel $i$
and channel $j$. Larger values indicate stronger estimated coupling between the
corresponding pair of channels.

All connectivity matrices were required to be square, finite, and symmetric
before being accepted by the pipeline. Symmetry is expected because the
connectivity measures used in this thesis are pairwise undirected measures:
the coupling between channel $i$ and channel $j$ is the same as the coupling
between channel $j$ and channel $i$. Therefore, each matrix represents an
undirected weighted network.

The diagonal entries of the connectivity matrices were set to one during
connectivity estimation, since each channel is perfectly connected with itself.
However, self-connections are not meaningful for graph-theoretical topology.
Consequently, self-loops were removed before graph metrics were computed by
setting the diagonal entries to zero in the graph construction stage.

The range of matrix values depended on the connectivity method. Coherence and
wPLI produced non-negative values in the interval $[0,1]$. AEC was based on
correlation coefficients and could therefore produce values in the interval
$[-1,1]$. Negative AEC values were retained when computing mean connectivity,
but negative values were not used as topology weights during graph-metric
calculation, as described in the graph feature extraction section.

\subsection{Channel-to-Node Mapping}
\label{subsec:channel-to-node-mapping}

A fixed channel-to-node mapping was used for all adjacency matrices. This was
necessary because graph metrics are only comparable across seizures and
subjects if the same node index always refers to the same EEG derivation.
The CHB-MIT dataset contains recordings with montage variability, and previous
CHB-MIT studies have therefore restricted analysis to fixed channel sets in
order to reduce heterogeneity across recordings \cite{Tsiouris2018}. In the
present thesis, a stricter screening strategy was used: recordings were retained
only when the complete predefined \num{22}-channel bipolar montage was
available. This preserved a larger sensor-level network while ensuring that the
node definitions remained identical across subjects, seizures, temporal
windows, frequency bands, and connectivity methods.

The channel order was inherited from the standardized preprocessing montage and
was kept fixed during segmentation, connectivity estimation, and graph feature
extraction. For readability, \cref{tab:channel-node-mapping} reports the
mapping using one-based node numbering. Internally, the computational
implementation used zero-based Python indices, but the ordering was identical.

\begin{table}[H]
    \centering
    \begin{tabular}{clcl}
        \toprule
        Node & Channel & Node & Channel \\
        \midrule
        1  & \texttt{FP1-F7}   & 12 & \texttt{P4-O2} \\
        2  & \texttt{F7-T7}    & 13 & \texttt{FP2-F8} \\
        3  & \texttt{T7-P7}    & 14 & \texttt{F8-T8} \\
        4  & \texttt{P7-O1}    & 15 & \texttt{T8-P8} \\
        5  & \texttt{FP1-F3}   & 16 & \texttt{P8-O2} \\
        6  & \texttt{F3-C3}    & 17 & \texttt{FZ-CZ} \\
        7  & \texttt{C3-P3}    & 18 & \texttt{CZ-PZ} \\
        8  & \texttt{P3-O1}    & 19 & \texttt{P7-T7} \\
        9  & \texttt{FP2-F4}   & 20 & \texttt{T7-FT9} \\
        10 & \texttt{F4-C4}    & 21 & \texttt{FT9-FT10} \\
        11 & \texttt{C4-P4}    & 22 & \texttt{FT10-T8} \\
        \bottomrule
    \end{tabular}
    \caption{Fixed channel-to-node mapping used for all functional connectivity
    matrices and graph-theoretical analyses.}
    \label{tab:channel-node-mapping}
\end{table}

Because this mapping was fixed, the same matrix position always represented the
same pair of bipolar EEG channels. For example, the value at position
$(1,2)$ represented the connection between \texttt{FP1-F7} and \texttt{F7-T7}
in every matrix, regardless of subject, seizure, window, frequency band, or
connectivity method.

\subsection{Frequency- and Window-Specific Networks}
\label{subsec:frequency-window-specific-networks}

The network construction procedure was repeated independently for each temporal
window and frequency band. Therefore, the analysis did not produce a single
network per seizure, but rather a structured collection of networks describing
the same seizure from different temporal, spectral, and methodological
perspectives.

For each seizure, the possible temporal windows were Baseline, T0, T1, and T2.
For each window, connectivity was computed in the delta, theta, alpha1,
alpha2, beta1, and beta2 bands. For each frequency band, three connectivity
methods were applied: coherence, wPLI, and AEC. A single adjacency matrix can
therefore be identified by the tuple

\[
    (\text{subject},\ \text{seizure},\ \text{window},\ \text{band},\
    \text{method}).
\]

This organization preserved the full experimental structure of the analysis.
It allowed graph metrics to be compared across preictal windows, across
frequency bands, and across connectivity methods while retaining the identity
of the subject/case folder and seizure from which each network was derived.

The resulting adjacency matrices were saved and subsequently passed to the
graph feature extraction stage. At that stage, the matrices were cleaned,
thresholded, converted into graph objects, and summarized using global
network-level metrics.

\section{Graph-Theoretical Metrics}
\label{sec:graph-theoretical-metrics}

Before computing topology-based graph metrics, each matrix was prepared in a
standardized way. First, the matrix was symmetrized and the diagonal was set to
zero, removing self-loops. For topology-based metrics, negative connectivity values were also set to zero.
This was done because the graph metrics used in this thesis interpret edge
weights as non-negative connection strengths. In particular, shortest-path
metrics converted connection strength into distance, so negative weights could
not be interpreted as valid positive distances in the same way as positive
connectivity values \cite{Rubinov2010,Rubinov2011}. However, negative values
were preserved when computing mean connectivity, so that the average
connectivity measure still reflected the original signed connectivity values
when applicable.

A proportional threshold was then applied to the cleaned matrix. Specifically,
the strongest positive edges were retained up to a target density of
\SI{20}{\percent}. Proportional thresholding was used to impose the same
intended graph density across subjects, seizures, temporal windows, frequency
bands, and connectivity methods. This reduced the influence of weak connections
that may be more susceptible to noise while making graph metrics more
comparable across matrices \cite{VanWijk2010,Garrison2015,Adamovich2022}.

The \SI{20}{\percent} density was selected a priori as a compromise between
sparsity and network preservation. A lower threshold would retain fewer edges
and could increase the risk of disconnected or unstable graphs, whereas a much
higher threshold would retain more weak connections and reduce the effect of
network sparsification. Since there is no universally accepted density threshold
for EEG functional connectivity graphs, the threshold should be interpreted as a
fixed methodological choice rather than an empirically optimal value
\cite{Garrison2015,Adamovich2022}. If a matrix contained fewer positive edges
than required by the target density, all available positive edges were retained
and the realized density was lower than the target.

In the resulting graph, nodes corresponded to EEG channels and edges
corresponded to retained functional connectivity values. Edge weights
represented connectivity strength. For metrics based on shortest paths,
connectivity strength was converted into distance using the inverse of the
edge weight, so that stronger functional connections corresponded to shorter
graph distances.

\subsection{Mean Connectivity}
\label{subsec:mean-connectivity}

Mean connectivity was computed directly from the original, unthresholded
connectivity matrix after symmetrization and removal of self-connections. It was
defined as the average of all off-diagonal upper-triangular matrix entries:

\[
    \bar{A}
    =
    \frac{1}{N(N-1)/2}
    \sum_{i<j} A_{ij},
\]

where $N$ is the number of nodes and $A_{ij}$ is the connectivity value between
channels $i$ and $j$.

This metric describes the overall level of functional coupling in a given
network \cite{Friston1994,Rubinov2010}. Unlike the topology-based graph metrics, mean connectivity was not
computed after thresholding. This was done to preserve information about the
average connectivity strength across all channel pairs. For AEC matrices,
negative correlations could therefore contribute to the mean connectivity
value.

\subsection{Density}
\label{subsec:density}

Graph density describes the proportion of possible edges that are present in
the thresholded graph \cite{Rubinov2010,VanWijk2010}.. For an undirected graph with $N$ nodes and $E$ retained
edges, density is defined as

\[
    D =
    \frac{E}{N(N-1)/2}.
\]

In this thesis, a proportional threshold targeting a density of
\SI{20}{\percent} was applied before topology metrics were computed. Therefore,
the density metric also served as a diagnostic measure confirming how many
edges were retained after thresholding. In most cases, the realized density was
expected to match the target density. If too few positive edges were available,
the realized density could be lower.

\subsection{Clustering Coefficient}
\label{subsec:clustering-coefficient}

The clustering coefficient measures the tendency of nodes to form locally
interconnected groups \cite{Watts1998,Onnela2005,Rubinov2010}. In the context of EEG functional networks, a high
clustering coefficient indicates that channels connected to the same channel
also tend to be connected to each other, suggesting stronger local network
organization.

For each thresholded weighted graph, the average weighted clustering
coefficient was computed across all nodes:

\[
    C =
    \frac{1}{N}
    \sum_{i=1}^{N} C_i,
\]

where $C_i$ is the weighted clustering coefficient of node $i$. The use of
weighted clustering allowed the metric to account not only for the presence of
connections but also for their connectivity strength.

\subsection{Global Efficiency}
\label{subsec:global-efficiency}

Global efficiency measures how efficiently information can be exchanged across
the whole network \cite{Latora2001,Rubinov2010}. It is based on the inverse of the shortest-path distance
between all pairs of nodes:

\[
    E_{\mathrm{glob}}
    =
    \frac{1}{N(N-1)}
    \sum_{i \neq j}
    \frac{1}{d_{ij}},
\]

where $d_{ij}$ is the shortest-path distance between nodes $i$ and $j$.

In this thesis, shortest-path distances were computed using inverse
connectivity strength as the edge distance. Therefore, stronger functional
connections corresponded to shorter graph distances. If two nodes were
disconnected, their contribution to global efficiency was set to zero. Higher
global efficiency indicates that the network has shorter paths between nodes
and may therefore support more integrated communication across the EEG channel
network.

\subsection{Characteristic Path Length}
\label{subsec:characteristic-path-length}

Characteristic path length measures the average shortest-path distance between
nodes in a network \cite{Watts1998,Rubinov2010}. It provides a complementary view to global efficiency:
whereas global efficiency increases when nodes are more easily reachable,
characteristic path length decreases when the network becomes more integrated.

For a connected graph, characteristic path length is defined as

\[
    L =
    \frac{1}{N(N-1)}
    \sum_{i \neq j}
    d_{ij},
\]

where $d_{ij}$ is the shortest-path distance between nodes $i$ and $j$.

Because thresholding may produce disconnected graphs, characteristic path
length was computed on the largest connected component of each graph. This
avoids undefined path lengths between disconnected node pairs. As with global
efficiency, edge distances were defined as the inverse of connectivity strength.

\subsection{Small-Worldness}
\label{subsec:small-worldness}

Small-worldness quantifies whether a network combines high local clustering
with short global path lengths compared with random reference networks \cite{Watts1998,Humphries2008,Rubinov2010}. This is
relevant for brain network analysis because many functional brain networks have
been described as balancing local specialization and global integration.

Small-worldness was computed as

\[
    \sigma =
    \frac{C / C_{\mathrm{rand}}}{L / L_{\mathrm{rand}}},
\]

where $C$ and $L$ are the clustering coefficient and characteristic path length
of the observed graph, and $C_{\mathrm{rand}}$ and $L_{\mathrm{rand}}$ are the
corresponding mean values from randomized comparison graphs.

For each observed graph, random comparison graphs were generated by preserving
the observed degree sequence through double-edge swaps when possible. The
observed edge-weight distribution was then shuffled onto each randomized
topology. In the final pipeline, \num{20} random graphs were used for each
small-worldness estimate, with a fixed random seed to ensure reproducibility.
Values greater than one indicate stronger small-world organization relative to
the random comparison graphs.

\subsection{Modularity}
\label{subsec:modularity}

Modularity measures the extent to which a network can be divided into
communities or modules \cite{Newman2006,Rubinov2010}. A module is a group of nodes that are more strongly
connected to each other than to the rest of the network. In EEG functional
networks, higher modularity may indicate stronger segregation of the network
into relatively independent channel groups.

Modularity can be written conceptually as

\[
    Q =
    \frac{1}{2m}
    \sum_{i,j}
    \left[
    A_{ij}
    -
    \frac{k_i k_j}{2m}
    \right]
    \delta(c_i,c_j),
\]

where $A_{ij}$ is the edge weight between nodes $i$ and $j$, $k_i$ and $k_j$
are node strengths or degrees, $m$ is the total edge weight, and
$\delta(c_i,c_j)$ indicates whether nodes $i$ and $j$ belong to the same
community.

In the implemented pipeline, modularity was computed using weighted community
detection. When available, the Louvain community detection method was used;
otherwise, a greedy modularity-based method was used as a fallback \cite{Blondel2008}. The same
random seed was used to make the modularity calculation reproducible when
stochastic community detection was applied.

\subsection{Mean Betweenness Centrality}
\label{subsec:mean-betweenness-centrality}

Betweenness centrality measures how often a node lies on shortest paths between
other nodes \cite{Freeman1977,Rubinov2010}. Nodes with high betweenness may act as bridges between different
parts of the network. Betweenness centrality has also been used in EEG connectivity studies of drug-resistant epilepsy and neuromodulation \cite{Lanzone2022}. 
In the present analysis, betweenness centrality was
computed for each node and then averaged across all nodes to obtain a single
network-level metric.

For a node $v$, betweenness centrality is defined as

\[
    BC(v)
    =
    \sum_{s \neq v \neq t}
    \frac{\sigma_{st}(v)}{\sigma_{st}},
\]

where $\sigma_{st}$ is the number of shortest paths between nodes $s$ and $t$,
and $\sigma_{st}(v)$ is the number of those paths that pass through node $v$.

Shortest paths were computed using inverse connectivity strength as distance,
so stronger edges corresponded to shorter paths. The mean betweenness
centrality was then computed as

\[
    \overline{BC}
    =
    \frac{1}{N}
    \sum_{v=1}^{N} BC(v).
\]

This produced a global summary of how strongly the network relied on central
bridge-like nodes.

\subsection{Assortativity}
\label{subsec:assortativity}

Assortativity measures the tendency of nodes to connect to other nodes with
similar degree \cite{Newman2002,Rubinov2010}. A positive assortativity coefficient indicates that highly
connected nodes tend to connect to other highly connected nodes, whereas a
negative coefficient indicates that highly connected nodes tend to connect to
less connected nodes.

Degree assortativity was computed for each thresholded graph. The coefficient
can be interpreted as a correlation between the degrees of nodes at either end
of an edge. Values close to one indicate assortative mixing, values close to
zero indicate no clear degree-based mixing pattern, and values below zero
indicate disassortative organization.

In the context of this thesis, assortativity was included as a global measure
of network organization complementary to efficiency, clustering, and modularity.
When assortativity was undefined, for example because the graph structure did
not contain sufficient degree variability, the value was recorded as missing.

\section{Output of the Connectivity and Graph Feature Pipeline}
\label{sec:connectivity-graph-pipeline-output}

The connectivity and graph feature pipeline produced the network-level dataset
used for statistical analysis. The final segmentation stage contained
\num{632} saved temporal windows. Each window was analyzed across \num{6}
frequency bands and \num{3} connectivity methods, giving

\[
    632 \times 6 \times 3 = \num{11376}
\]

expected connectivity matrices. No matrices failed validation, so the final
connectivity output consisted of \num{11376} valid \num{22} $\times$ \num{22}
matrices.

Each matrix was then summarized by graph-theoretical features. The pipeline
retained the subject/case identifier, seizure identifier, temporal window,
frequency band, and connectivity method for every matrix, so each graph feature
row remained traceable to its source EEG segment and connectivity definition.

The graph feature extraction stage produced one row per connectivity matrix,
resulting in \num{11376} graph-level observations. The extracted variables were
mean connectivity, density, clustering coefficient, global efficiency,
characteristic path length, small-worldness, modularity, mean betweenness
centrality, and assortativity. Density was retained as a diagnostic measure,
whereas the remaining graph metrics were used as outcome variables in the
statistical analysis.

Additional diagnostic variables were stored for quality control and
reproducibility. These included the threshold density, realized edge count,
realized graph density, number of connected components, size of the largest
connected component, number of random graphs used for small-worldness
estimation, random seed, and software version information.

Thus, the final output of this stage was a seizure-window-level graph feature
dataset. Each row represented one functional network defined by a specific
subject/case identifier, seizure, temporal window, frequency band, and
connectivity method. This dataset provided the input for the statistical
analysis in the following chapter, where graph metrics were compared across
Baseline, T0, T1, and T2.
